\section{Análisis comparativo de VE-SDE y VP-SDE}

\subsection{Diferencias clave entre dos SDE}

Antes de profundizar en Probability Flow ODE, comparemos sistemáticamente las diferencias entre VE-SDE y VP-SDE; esto facilita la comprensión de los compromisos de diseño en distintos modelos de difusión.

\begin{table}[H]
\centering
\begin{tabular}{lcc}
\toprule
\textbf{Característica} & \textbf{VE-SDE (SMLD)} & \textbf{VP-SDE (DDPM)} \\
\midrule
Término de deriva $f(\mathbf{x},t)$ & $0$ & $-\frac{1}{2}\beta(t)\mathbf{x}$ \\
Coeficiente de difusión $g(t)$ & $\sqrt{\mathrm{d}[\sigma^2]/\mathrm{d}t}$ & $\sqrt{\beta(t)}$ \\
Atenuación de la media & No ($\alpha_t = 1$) & Sí ($\alpha_t < 1$) \\
Comportamiento de la varianza & Crecimiento monótono (exploding) & Mantenimiento constante (preserving) \\
Distribución de estado final & $\mathcal{N}(0, \sigma_T^2 \mathbf{I})$ & $\mathcal{N}(0, \mathbf{I})$ \\
Modo de atenuación del SNR & $1/\sigma^2(t)$ & $\alpha_t^2/(1-\alpha_t^2)$ \\
\bottomrule
\end{tabular}
\caption{Comparativa sistemática de VE-SDE y VP-SDE}
\end{table}

\begin{knowledgebox}{Intuición física}
La imagen física del \textbf{VE-SDE} es: cada punto de datos permanece "estacionario", pero la niebla de ruido que lo rodea se vuelve cada vez más densa; la varianza aumenta continuamente hasta ocultar por completo la estructura de los datos. La imagen física del \textbf{VP-SDE} corresponde a un proceso Ornstein-Uhlenbeck: los puntos de datos son atraídos simultáneamente hacia el origen por una fuerza elástica (atenuación de la media) y experimentan perturbaciones aleatorias (inyección de ruido); ambos factores se equilibran para mantener la varianza total constante.
\end{knowledgebox}

\subsection{Comparativa del rendimiento experimental}

Los experimentos de Song et al. (2021) muestran que, sobre el conjunto de datos CIFAR-10:
\begin{itemize}
    \item Al emplear muestreo DDPM con 1000 pasos, el FID (Frechet Inception Distance, donde menor es mejor) oscila alrededor de 3--4;
    \item Con DDIM, se alcanza un FID similar utilizando apenas unos 50 pasos;
    \item Mediante el uso de un solver SDE mejorado (por ejemplo, basado en métodos Predictor-Corrector), una calidad aceptable se logra con tan solo 10--15 pasos.
\end{itemize}

\begin{warningbox}{El FID no es la única métrica}
El FID mide simultáneamente la calidad generativa y la diversidad, pero no refleja completamente la percepción humana de la calidad de las imágenes. En aplicaciones prácticas, también deben considerarse la velocidad de generación, la cobertura de modos (diversity) y el desempeño en tareas posteriores. VE-SDE y VP-SDE presentan ventajas y desventajas relativas según la configuración utilizada.
\end{warningbox}

\subsection{Resumen de la sección}

VE-SDE y VP-SDE representan dos estrategias de ruido distintas: la primera implica difusión pura sin deriva, mientras que la segunda equilibra deriva y difusión. Comprender sus diferencias ayuda a tomar decisiones de diseño razonables en aplicaciones prácticas.
